\chapter{बड़ी संख्याओं का गुणा}\label{ch:g4:mult}

पिछले साल हमने एक अंक से गुणा किया था (\cref{ch:g3:mult}); यह
अध्याय दो अंकों की संख्याओं से गुणा करता है --- और खानों वाले
तरीक़े की रहस्यमयी ``खिसकी हुई पंक्ति'' समझाता है, जो दहाई और
इकाई में तोड़ने के सिवा कुछ नहीं है।

\section{दहाइयों से गुणा}

\begin{proposition}[10, 20, 300 आदि से गुणा]\label{prop:g4:mult:tens}
$10$ से गुणा करने पर एक \omterm{def:g1:counting:zero}{शून्य} जुड़ता है, $100$ से दो। $20$ से
गुणा करने के लिए पहले $2$ से गुणा करो, फिर $10$ से:
\[
36 \times 20 = 36 \times 2 \times 10 = 72 \times 10 = 720 .
\]
इसी तरह $14 \times 300 = 14 \times 3 \times 100 = 4\,200$।
\end{proposition}
\admitted

\section{दो अंकों वाला खानों का तरीक़ा}

\begin{example}[खिसकी हुई पंक्ति क्यों?]\label{ex:g4:mult:why}
$23 \times 12$ निकालने के लिए $12$ को $10 + 2$ में तोड़ो:
\[
23 \times 12 = 23 \times 2 + 23 \times 10 = 46 + 230 = 276 .
\]
खानों वाला तरीक़ा ठीक यही दो \omterm{def:g2:mult:def}{गुणनफल} एक के नीचे एक लिखता है ---
दूसरा बाईं ओर खिसका हुआ, क्योंकि वह \emph{दहाइयों} की संख्या
है:
\[
\begin{array}{r}
2\,3 \\
\times\ 1\,2 \\
\hline
4\,6 \\
2\,3\,\cdot \\
\hline
2\,7\,6
\end{array}
\]
(बिंदु खिसकी हुई पंक्ति का इकाई स्थान दिखाता है, जिसे अक्सर
ख़ाली छोड़ दिया जाता है)।
\end{example}

\begin{omfigure}
\begin{tikzpicture}[scale=0.42]
  \fill[omDef!18] (0,0) rectangle (20,10);
  \fill[omThm!18] (0,10) rectangle (20,12);
  \fill[omDef!34] (20,0) rectangle (23,10);
  \fill[omThm!34] (20,10) rectangle (23,12);
  \draw[thick] (0,0) rectangle (23,12);
  \draw[thick] (20,0) -- (20,12);
  \draw[thick] (0,10) -- (23,10);
  \node[font=\small] at (10,5) {$20 \times 10 = 200$};
  \node[font=\small] at (10,11) {$20 \times 2 = 40$};
  \node[font=\small, rotate=90] at (21.5,5) {$3 \times 10 = 30$};
  \node[font=\small] at (21.5,11) {$6$};
  \node[below, font=\small] at (10,0) {$20$};
  \node[below, font=\small] at (21.5,0) {$3$};
  \node[left, font=\small] at (0,5) {$10$};
  \node[left, font=\small] at (0,11) {$2$};
\end{tikzpicture}

{\small $23 \times 12$ का आयत-चित्र: चार आसान \omterm{def:g2:mult:def}{गुणनफल},
$200 + 40 + 30 + 6 = 276$। खानों वाला तरीक़ा इन्हें दो
पंक्तियों में समेट देता है: $46$ (``$\times 2$'' वाली पंक्ति)
और $230$ (``$\times 10$'' वाली पंक्ति)।}
\end{omfigure}

\begin{method}[दो अंकों की संख्या से खानों में गुणा]\label{met:g4:mult:column}
$457 \times 36$ निकालने के लिए:
\begin{enumerate}
  \item पहली पंक्ति: $457 \times 6 = 2\,742$;
  \item दूसरी पंक्ति: $457 \times 3$ \emph{दहाई} $= 1\,371$
        दहाई --- $1\,371$ को एक स्थान बाईं ओर खिसकाकर लिखो;
  \item दोनों पंक्तियाँ जोड़ो: $2\,742 + 13\,710 = 16\,452$;
  \item जाँच के लिए अंदाज़ा: $457 \times 36 \approx 500 \times 36
        = 18\,000$? इससे बेहतर: $450 \times 36$ लगभग $16\,000$
        है --- भरोसेमंद।
\end{enumerate}
\[
\begin{array}{r}
4\,5\,7 \\
\times\ \ \,3\,6 \\
\hline
2\,7\,4\,2 \\
1\,3\,7\,1\,\cdot \\
\hline
1\,6\,4\,5\,2
\end{array}
\]
\end{method}

\begin{example}[मानसिक तरकीबें]\label{ex:g4:mult:tricks}
\begin{itemize}
  \item \emph{मिलनसार ढंग से तोड़ो}: $18 \times 11 = 18 \times 10
        + 18 = 198$।
  \item \emph{सौ के पड़ोसी का उपयोग}: $25 \times 99 = 25 \times 100 -
        25 = 2\,475$।
  \item \emph{एक को दुगुना, दूसरे को \omterm{def:g3:division:half}{आधा}}: $16 \times 35 = 8 \times 70 =
        560$ (एक गुणक को \omterm{def:g3:division:half}{आधा} और दूसरे को दुगुना करने से
        \omterm{def:g2:mult:def}{गुणनफल} वही रहता है)।
\end{itemize}
\end{example}

\section{समस्याओं में गुणा}

\begin{example}\label{ex:g4:mult:problem}
एक विद्यालय $145$-$145$ काग़ज़ों के $28$ डिब्बे मँगाता है।
\begin{enumerate}
  \item संक्रिया: $145 \times 28$।
  \item अंदाज़ा: $150 \times 30 = 4\,500$, दोनों ओर से थोड़ा
        ज़्यादा।
  \item खाने: $145 \times 8 = 1\,160$; $145 \times 2$ दहाई
        $= 2\,900$; कुल $4\,060$।
  \item वाक्य: विद्यालय को $4\,060$ काग़ज़ मिलते हैं।
\end{enumerate}
\end{example}

\section{अभ्यास}

\begin{exercise}[$\star$]\label{exo:g4:mult:1}
मन-ही-मन निकालो: $47 \times 10$; \quad $23 \times 100$; \quad
$36 \times 20$; \quad $15 \times 300$; \quad $42 \times 50$।
\end{exercise}

\begin{exercise}[$\star$]\label{exo:g4:mult:2}
$34 \times 21$ को \cref{ex:g4:mult:why} की तरह तोड़कर निकालो
($34 \times 21 = 34 \times 20 + 34$), फिर खानों में जाँचो।
\end{exercise}

\begin{exercise}[$\star$]\label{exo:g4:mult:3}
खानों में निकालो: $63 \times 24$; \quad $85 \times 47$; \quad
$126 \times 53$।
\end{exercise}

\begin{exercise}[$\star$]\label{exo:g4:mult:4}
पहले अंदाज़ा लगाकर खानों में निकालो: $208 \times 34$; \quad
$517 \times 62$।
\end{exercise}

\begin{exercise}[$\star$]\label{exo:g4:mult:5}
\cref{ex:g4:mult:tricks} की किसी तरकीब से निकालो:
$45 \times 11$; \quad $32 \times 99$; \quad $24 \times 45$
(दुगुना और \omterm{def:g3:division:half}{आधा} --- चाहो तो दो बार)।
\end{exercise}

\begin{exercise}[$\star$]\label{exo:g4:mult:6}
एक सिनेमाघर में $18$-$18$ सीटों की $26$ पंक्तियाँ हैं। कुल कितनी सीटें?
\end{exercise}

\begin{exercise}[$\star$]\label{exo:g4:mult:7}
एक ट्रक $48$-$48$ पेटियों के $32$ फट्टे ढोता है। कुल कितनी
पेटियाँ? पहले अंदाज़ा दो, फिर ठीक उत्तर।
\end{exercise}

\begin{exercise}[$\star$]\label{exo:g4:mult:8}
ग़लती ढूँढ़ो: अना $57 \times 23$ ऐसे निकालती है:
$57 \times 3 = 171$ और $57 \times 2 = 114$, फिर
$171 + 114 = 285$। सही उत्तर $1\,311$ है। वह क्या भूल
गई?
\end{exercise}

\begin{exercise}[$\star\star$]\label{exo:g4:mult:9}
एक माली $35$-$35$ ट्यूलिप की $14$ पंक्तियाँ और $28$-$28$ नरगिस
की $12$ पंक्तियाँ लगाता है। कुल कितने फूल? (दो \omterm{def:g2:mult:def}{गुणनफल}, फिर एक \omterm{def:g1:addition:def}{योगफल}।)
\end{exercise}

\begin{exercise}[$\star\star$]\label{exo:g4:mult:10}
एक नानबाई रोज़ $67$ kg आटा इस्तेमाल करता है। $30$ दिनों के
महीने में यह लगभग कितना आटा हुआ --- गोल की हुई संख्याओं से
अंदाज़ा लगाओ, फिर ठीक-ठीक निकालो।
\end{exercise}

\begin{exercise}[$\star\star$]\label{exo:g4:mult:11}
इस अध्याय जैसे आयत-चित्र से समझाओ कि $15 \times 12$ को
$15 \times 10 + 15 \times 2$ के रूप में क्यों निकाला जा सकता है।
फिर $15 \times 12$ को तोड़ने का एक और चलने वाला तरीक़ा सोचो।
\end{exercise}
